% =============================================================================
% Relatorio 2, item D: estudo de aproveitamento tecnico (projeto Riachuelo).
% Benchmarking por linhas de similaridade (L1 a L8), solucoes descartadas e
% propriedade intelectual (familias conferidas no Google Patents em 25/09/2026).
% Tabela: tab:r2-aproveitamento (longtable). Mapa de posicionamento e vantagens
% comerciais ficam no estudo de diferenciacao (r2/5-diferenciacao.tex).
% =============================================================================

\section{ESTUDO DE APROVEITAMENTO TÉCNICO}
\label{sec:r2-aproveitamento}

O equilíbrio econômico da Seção~\ref{sec:r2-valor} depende de cumprir o custo-meta e de reduzir o risco técnico do desenvolvimento, e as duas coisas ficam mais fáceis quando o sistema reaproveita soluções já testadas em outros produtos.

\subsection{Método}
\label{sub:r2-aprov-metodo}

O estudo de aproveitamento técnico procurou, em produtos análogos, patentes e casos de mercado, soluções já testadas que o sistema Tag\&Go pudesse incorporar em vez de desenvolver do zero. A orientação da disciplina apontou o benchmarking de produtos análogos e de patentes, inclusive no INPI, como a fonte mais útil de princípios de solução e recomendou, em funções como a recomendação de roupas, construir sobre bases existentes \cite{zancul2026conceitual}. A equipe seguiu o precedente de um grupo de anos anteriores da disciplina, o do projeto Carrinho de compras, e organizou a busca por linhas de similaridade: cada linha reúne referências que resolvem um mesmo problema técnico do Tag\&Go, ainda que em outro produto ou em outro setor.

Para cada linha, a equipe listou as referências, extraiu a ideia aproveitável e indicou onde ela entra no sistema, pela função da análise funcional (Tabela~\ref{tab:r2-funcoes}) e pelo subsistema ou componente da arquitetura (Tabela~\ref{tab:r2-ssc}). O critério de aproveitamento foi técnico: a ideia precisava realizar uma função do Tag\&Go, ser compatível com as especificações-meta e com o envelope herdado do Relatório 1 e não depender de recurso sem precedente comercial. As oito linhas resultantes, L1 a L8, cobrem a etiqueta e sua liberação (L1, L2, L3 e L5), o pagamento (L4), a loja (L6), o ciclo logístico (L7) e a Jornada Completa (L8), e estão consolidadas na Tabela~\ref{tab:r2-aproveitamento}. O mapa de posicionamento e as vantagens comerciais frente às mesmas referências pertencem ao estudo de diferenciação (Seção~\ref{sec:r2-diferenciacao}); o aproveitamento técnico trata apenas do que cada referência ensina sobre como realizar as funções. As ideias aproveitadas alimentaram os princípios de solução e a matriz morfológica da Seção~\ref{sec:r2-concepcao}.

\subsection{Linhas de similaridade}
\label{sub:r2-aprov-linhas}

A primeira linha, retenção mecânica (L1), partiu das etiquetas rígidas comerciais e da garra e do pino retirados de mini tags doadoras no protótipo do Relatório 1. Essas etiquetas usam uma embreagem de três esferas num copo cônico, em que a retenção cresce com a força de extração: quanto mais se puxa o pino, mais as esferas apertam a haste. Esse autotravamento dispensa energia para reter a peça, o que é decisivo para uma etiqueta que passa a maior parte do ciclo parada no estado TRAVADA. A equipe aproveitou a garra (\O\,12,4 $\times$ 10\,mm, curso de 1,2\,mm e mola estimada em 2,5\,N, valores a medir), o pino padrão do mercado, com cabeça de \O\,14\,mm e haste de \O\,2,0\,mm, e o elemento EAS da mini tag, que mantém a compatibilidade com os pórticos EAS já instalados nas lojas. O pino padrão preserva, além disso, o gesto de aplicação que já é conhecido de quem etiqueta peças, o que simplifica o cadastro na origem. A única mudança na garra é o carretel não ferromagnético, que impede a abertura pelo destacador comercial (especificação E11). A linha entra nas funções reter peça (F1) e sinalizar saída indevida (F12) e nos componentes C1.3 (garra), C1.4 (pino) e C1.10 (elemento EAS). A meta de resistência à extração de pelo menos 120\,N (E02) é da equipe e ainda será comparada, em ensaio, com a de uma etiqueta comercial.

A segunda linha, liberação autorizada (L2), reuniu dispositivos que só soltam a peça depois de uma verificação. A patente US 7.450.013 B2, da Checkpoint, expirada em 2025 e hoje em domínio público \cite{googlepatents2008us7450013}, descreve um destacamento autenticado que depende de condições controladas, como um ímã propositalmente forte e uma velocidade calibrada, e prevê uma posição sem energia que permite a liberação de emergência. A SuperTag 4, da Sensormatic, confere se o item pago corresponde ao da etiqueta antes de acionar o destacador motorizado do self-checkout e, segundo o fornecedor, só destaca com compra válida \cite{sensormatic2026supertag}. A patente US 11.417.186 B2, da Kohl's, automatiza a retirada e deixa as etiquetas caírem em caixas coletoras \cite{googlepatents2022us11417186}, e as travas LIVE Locks, da InVue, abertas pelo celular, registram cada abertura por loja, departamento e usuário \cite{invue2026locks}. Dessas referências a equipe tirou quatro ideias. A primeira é que toda etiqueta precisa de um caminho de liberação independente da própria bateria: no Tag\&Go, o Liberador de Balcão energiza a etiqueta pelos contatos da face superior e entrega a autorização por eles. A segunda é a conferência entre etiqueta e item, feita pela associação entre o \texttt{tag\_id} e o EPC da peça na Estação de Cadastro, de modo que o preço venha sempre do EPC. A terceira é a queda por gravidade numa caixa trancada, que originou o Coletor de Etiquetas com boca antirretorno. A quarta é a trilha de auditoria, convertida na meta de registrar 100\% das emissões de token e das liberações (S10). A linha entra nas funções validar autorização (F5), oferecer liberação assistida (F16), recolher etiqueta (F17) e associar etiqueta à peça (F21) e nos componentes C1.5 (atuador), C4.1 (Coletor de Etiquetas) e C4.2 (Liberador de Balcão).

A terceira linha, identificação e leitura pelo celular (L3), tratou do caminho entre a peça e o canal de compra. A Apple documenta que o iPhone XS e os modelos posteriores leem NFC em segundo plano e, quando o primeiro registro NDEF traz um universal link, abrem o app ou o App Clip associado; sem app associado, abrem o Safari \cite{apple2026corenfc}. O App Clip Code combina código visual e NFC, mas a versão com NFC exige uma etiqueta NFC Type 5 de pelo menos 35\,mm \cite{apple2026appcliphig}, maior que a área disponível sob o QR de 28\,mm. O NTAG 424 DNA, da NXP, gera com AES-128 uma mensagem autenticada diferente a cada toque, o que impede copiar a URL \cite{nxp2026ntag424}. A equipe aproveitou a combinação de QR impresso e NFC com a mesma URL, que contém apenas o identificador público e estável da etiqueta; a peça atual vem da Plataforma Tag\&Go a cada leitura. A cópia da URL, por si só, não libera nenhuma peça, porque a liberação depende do token verificado na etiqueta (Seção~\ref{sec:r2-arquitetura}). Ainda assim, a equipe registrou o NTAG 424 DNA como evolução, caso a clonagem de URLs se mostre um problema na operação. A linha entra na função identificar peça (F2) e no componente C1.9 (identificação).

A quarta linha, pagamento sem app (L4), buscou formas de pagar sem instalar um aplicativo. Na Flying Tiger, a MishiPay oferece scan and go num web app aberto por QR, sem baixar app, e relata cesta 35\% maior, segundo o fornecedor \cite{mishipay2026flyingtiger}. Na apresentação dos App Clips na WWDC20, a Apple mostrou um toque numa etiqueta NFC que abre o App Clip na página de pedido e uma compra concluída com Apple Pay \cite{apple2020wwdcappclips}. Na web, a biblioteca JavaScript da Google Pay API permite pagar sem instalar app, desde que haja um processador de pagamento (PSP) \cite{google2026paywebvisao}. O Pix completa o conjunto: no recorte brasileiro de uma pesquisa da Adyen, com amostra não informada, 60\% dos consumidores disseram desistir da compra quando não podem pagar como querem \cite{tiinside2026adyen}. A equipe aproveitou esses elementos na Compra Rápida (Subseção~\ref{sub:r2-caminhos}): o App Clip Tag\&Go no iPhone e a Página de Compra Rápida no Android pagam pela carteira digital, com um toque no botão da carteira e confirmação por Face ID, Touch ID ou senha na folha do sistema, e oferecem o Pix copia e cola como alternativa. A linha entra na função processar pagamento (F3) e nos componentes C3.1 (App Clip Tag\&Go) e C3.2 (Página de Compra Rápida).

A quinta linha, energia e recarga (L5), comparou dispositivos que dependem de bateria ou de energia externa. As etiquetas eletrônicas de prateleira (ESL) da Vusion e da SOLUM são o precedente mais próximo de um dispositivo com bateria e rádio instalado em grande número no salão de venda \cite{vusion2026esl,solum2026esl}, e as travas da InVue anunciam bateria de dois anos \cite{invue2026locks}. A patente US 10.497.238 B2, da Sensormatic, descreve uma estação que energiza uma etiqueta sem bateria no momento da liberação \cite{googlepatents2019us10497238}. Os estudos anteriores da equipe mostraram, por um balanço de energia, que o gargalo de uma etiqueta sem bateria é a potência instantânea, e não a energia total: a RF ambiente, de cerca de 5\,$\mu$W, levaria cerca de 35\,h para acumular a energia de uma liberação com fio SMA; um campo NFC dedicado, de cerca de 6\,mW, cerca de 105\,s; e contatos de 5\,W, cerca de 1,5\,s. Dessas referências a equipe tirou três ideias: recarregar a célula por contatos, numa Bandeja de Recarga no CD, sem conectar cabos a cada etiqueta; manter a etiqueta em sono profundo e acordá-la só pelo botão (modo botão), o que dá autonomia de 448 dias com a célula de 150\,mAh, estimativa da equipe, contra a meta de 168 dias (E05); e fornecer energia externa pelo Liberador de Balcão quando a célula falhar. A linha entra nas funções armazenar energia (F7), repor energia (F8) e oferecer liberação assistida (F16) e nos componentes C1.7 (célula), C4.2 (Liberador de Balcão) e C5.3 (Bandeja de Recarga).

A sexta linha, localização e estoque (L6), comparou formas de informar ao cliente onde está a peça. O Store Mode da Zara mostra no mapa da loja onde a peça está exposta \cite{inditex2021storemode}; o app do Walmart marca o corredor num mapa em mais de 4.700 lojas \cite{walmart2019mapas}; e a Renner atualiza o estoque RFID a cada 60\,s \cite{infomoney2024renner}. Para maior precisão, o leitor de teto Impinj xArray localiza 85\% das etiquetas a até 1,5\,m em condição ideal, mas está em fim de vida e tem preço sob consulta \cite{atlasrfid2026xarray}, e uma solução comercial de Bluetooth com ângulo de chegada, citada pelo Bluetooth SIG, tem erro médio de 1 a 2\,m \cite{bluetooth2022direction}. A rapitag usou um algoritmo que achava pelo BLE o produto mais próximo do celular \cite{storesshops2019rapitag}, e as ESLs da Vusion têm LED de sete cores para guiar a separação e o cliente \cite{vusion2026esl}. Como localizar o produto na loja é uma função secundária acessória, a decisão foi de granularidade, entre o nível N0, da unidade, e o nível N3, da peça, e a escolha está na análise funcional (Tabela~\ref{tab:r2-granularidade}). A equipe aproveitou o nível N1, por setor, a partir do estoque RFID que a Riachuelo já opera e de um planograma digital mantido pela loja, como na Zara e no Walmart, e registrou o nível N2, por arara ou expositor, como evolução caso a Riachuelo instale leitores de teto. A ideia da rapitag não foi levada à etiqueta, porque o anúncio BLE contínuo esgotaria a célula antes do fim de um ciclo (Subseção~\ref{sub:r2-localizar}). A linha entra nas funções localizar produto na loja (F25) e informar disponibilidade (F26) e no serviço Jornada (C2.5).

A sétima linha, logística reversa (L7), estudou ativos retornáveis do varejo. A Sensormatic recircula etiquetas rígidas: as etiquetas e os pinos retirados no caixa são encaixotados, voltam aos CDs no retorno dos caminhões (\textit{back-haul}), são testados, limpos, separados, reembalados e enviados às fábricas de roupa; desde 2010, segundo o fornecedor, 16,5 bilhões de etiquetas foram enviadas e 13,7 bilhões recirculadas \cite{sensormatic2024recirculacao}, razão de cerca de 83\% calculada pela equipe. A Pact informa que suas etiquetas EAS reutilizáveis fazem no mínimo 8 a 10 viagens, com reuso de até 80\% \cite{pact2026etiquetas}, e a Inditex contratou a Tyco em 2014 para recircular etiquetas rígidas com RFID e AM em três a quatro centros, com apagamento do identificador \cite{rfidjournal2014inditex}. Nos cabides, a Braiform estima que reusar um cabide nove vezes reduz 79\% das emissões \cite{braiform2018cabides}, e a Mainetti opera o reuso desde meados dos anos 1980 \cite{mainetti2018cabides}. A Brambles, que opera paletes em pool, mostra que a taxa de perda do ativo retornável domina a economia do sistema \cite{brambles2025resultados}, e a Decathlon aplica o RFID na fábrica em 600 milhões de produtos por ano \cite{rfidnews2026decathlon}. A equipe aproveitou o Malote de Retorno no retorno do caminhão da frota própria, a Bancada de Inspeção no CD, o cadastro na origem, feito na Estação de Cadastro nos ramos fábrica e CD como no \textit{source tagging} da Decathlon, e um indicador de perda por loja, que acompanha a meta de retorno de 99\% (S01), acima da faixa de 80\% a 83\% observada no mercado. O Tag\&Go recircula uma etiqueta com bateria e chave criptográfica, o que acrescenta ao ciclo a recarga e as regras de transporte de baterias de lítio (Subseção~\ref{sub:r2-ciclo}); o efeito da taxa de retorno sobre a economia do sistema foi tratado na Seção~\ref{sec:r2-valor}. A linha entra nas funções F17 a F24 e no subsistema SS5 (ciclo logístico).

A oitava linha, jornada, fidelidade e dados (L8), reuniu referências para as funções da Jornada Completa. O Nike App at Retail reconhece o membro na loja, lê o produto para mostrar cores e tamanhos disponíveis, reserva peças para provar e permite pagar sem fila \cite{marketingdive2018nike}. A Sephora organiza a fidelidade em três níveis por gasto anual \cite{modernretail2024sephora}. O ``Complete the Look'' da Zalando aumentou em 25\% as interações com looks de criadores em dois meses \cite{zalando2021looks}, e a Stitch Fix combina filtragem colaborativa, similaridade visual e um stylist humano na seleção final \cite{stitchfix2026algoritmos}. A Renner usa algoritmos que comparam produtos por imagem e recomendação por transações \cite{infomoney2024renner}. Do lado da Riachuelo, o Ria Club já paga cashback de 12\% com o cartão Riachuelo e de 10\% nos demais meios \cite{midway2026riaclub}. O Apple Wallet aceita passes de fidelidade distribuídos pela web, por QR, por NFC ou por App Clip, e o Google Wallet emite cartões de fidelidade por link, ambos sem exigir o app do varejista \cite{apple2026walletpass,google2026walletfidelidade}. Seguindo a orientação de construir a recomendação sobre bases existentes \cite{zancul2026conceitual}, a equipe aproveitou quatro ideias: um modelo híbrido de recomendação (regras de coleção, filtragem colaborativa e similaridade visual, com looks de stylists) restrito ao que existe na loja no tamanho do perfil; o botão ``chamar vendedor'', que preserva o atendimento humano; o passe do Ria Club no Wallet, oferecido depois da Compra Rápida; e o cashback liberado só depois do registro de ``etiqueta devolvida'', que liga a fidelidade à meta de retorno. As demais funções da Jornada Completa, como a compra on-line do tamanho ausente e as métricas de sustentabilidade por categoria, seguem a Subseção~\ref{sub:r2-caminhos}. Essas funções continuam hipóteses, porque o survey do Relatório 1 não mediu interesse por fidelidade, recomendação ou sustentabilidade. A linha entra nas funções acionar atendimento humano (F15) e F27 a F31 e nos componentes C2.5 (Jornada), C3.3 (módulo Tag\&Go do app Riachuelo) e C3.4 (passe do Ria Club).

A partir das oito linhas, a equipe concluiu que quase todas as funções do Tag\&Go têm precedente comercial isolado: a retenção vem das etiquetas rígidas, a leitura e o pagamento sem app vêm das plataformas móveis, a recirculação vem da Sensormatic e da Pact, e a localização por setor vem da Zara e do Walmart. O desenvolvimento próprio fica concentrado na integração dessas soluções e em dois pontos para os quais a equipe não encontrou solução pronta para aproveitar: a verificação do token na própria etiqueta, com desafio gerado por ela, e a recirculação de uma etiqueta com bateria e chave criptográfica. A Tabela~\ref{tab:r2-aproveitamento} consolida as linhas, as ideias aproveitadas e o lugar de cada uma no sistema, com as funções da Tabela~\ref{tab:r2-funcoes} e os componentes da Tabela~\ref{tab:r2-ssc}.

\begingroup
\footnotesize
\setlength{\tabcolsep}{4pt}
\renewcommand{\arraystretch}{1.25}
\begin{longtable}{@{}
    >{\RaggedRight\arraybackslash}p{2.2cm}
    >{\RaggedRight\arraybackslash}p{5.3cm}
    >{\RaggedRight\arraybackslash}p{4.6cm}
    >{\RaggedRight\arraybackslash}p{2.5cm}@{}}
\caption{Linhas de similaridade e ideias aproveitadas} \label{tab:r2-aproveitamento} \\
\toprule
\textbf{Linha} & \textbf{Referências} & \textbf{Ideia aproveitada} & \textbf{Onde entra} \\
\midrule
\endfirsthead
\multicolumn{4}{c}{Tabela~\ref{tab:r2-aproveitamento} (continuação)} \\
\toprule
\textbf{Linha} & \textbf{Referências} & \textbf{Ideia aproveitada} & \textbf{Onde entra} \\
\midrule
\endhead
\midrule
\multicolumn{4}{r}{\textit{continua na próxima página}} \\
\endfoot
\bottomrule
\multicolumn{4}{c}{Fonte: Elaborado pelos autores (2026) com base nas fontes citadas no texto.} \\
\endlastfoot
L1 Retenção mecânica & etiquetas rígidas comerciais; garra e pino retirados de mini tags no protótipo do Relatório 1 & embreagem de 3 esferas com copo cônico, pino padrão do mercado e elemento EAS de mini tag: compatibilidade com a operação atual e com os pórticos & F1, F12; C1.3, C1.4, C1.10 \\
\addlinespace
L2 Liberação autorizada & US 7.450.013 B2 (Checkpoint, domínio público desde 2025): ímã propositalmente forte, velocidade calibrada e posição sem energia que permite liberação de emergência; SuperTag 4: conferência item-etiqueta; Kohl's US 11.417.186 B2: retirada automática com queda em caixas coletoras; InVue: trilha de auditoria & liberação de emergência pelo Liberador de Balcão; conferência etiqueta-EPC; Coletor de Etiquetas com queda por gravidade; registro de toda liberação (S10) & F5, F16, F17, F21; C1.5, C4.1, C4.2 \\
\addlinespace
L3 Identificação e leitura pelo celular & NFC com universal link e QR (Apple); App Clip Code; NTAG 424 DNA (mensagem autenticada a cada toque) & mesma URL no QR e no NFC; evolução para NFC com mensagem autenticada se houver clonagem de URL & F2; C1.9 \\
\addlinespace
L4 Pagamento sem app & MishiPay (web app por QR na Flying Tiger); Apple Pay em App Clip (WWDC20); Google Pay API na web; Pix & Compra Rápida por App Clip e página; carteira com um toque e confirmação por Face ID, Touch ID ou senha; Pix copia e cola como alternativa & F3; C3.1, C3.2 \\
\addlinespace
L5 Energia e recarga & ESLs com bateria e rádio (Vusion, SOLUM); InVue (bateria de 2 anos); US 10.497.238 B2 (estação que energiza etiqueta sem bateria); balanço de energia dos estudos anteriores da equipe & recarga por contatos em bandeja no CD; modo de sono profundo; energia externa no Liberador de Balcão & F7, F8, F16; C1.7, C4.2, C5.3 \\
\addlinespace
L6 Localização e estoque & Zara Store Mode (planta); Walmart (corredor); Renner (estoque RFID a cada 60\,s); xArray; BLE com ângulo de chegada; rapitag (produto mais próximo pelo BLE); ESLs com LED & N1 por estoque RFID e planograma; N2 como evolução com leitores de teto & F25, F26; C2.5 \\
\addlinespace
L7 Logística reversa & Sensormatic (\textit{back-haul}, teste, limpeza, envio às fábricas); Pact (8 a 10 viagens; reuso de até 80\%); Inditex e Tyco (3 a 4 centros; apagamento do identificador); cabides Braiform e Mainetti; Brambles (a perda domina o custo); Decathlon (RFID na fábrica) & Malote de Retorno no retorno do caminhão; Bancada de Inspeção; cadastro na origem; indicador de perda por loja & F17 a F24; SS5 \\
\addlinespace
L8 Jornada, fidelidade e dados & Nike App at Retail; Sephora (níveis por gasto); Zalando (Complete the Look); Stitch Fix (algoritmo e stylist); Renner (CRM e IA); Ria Club; passes do Wallet & recomendação híbrida restrita ao estoque; botão ``chamar vendedor''; passe do Ria Club sem app; cashback após ``etiqueta devolvida'' & F15, F27 a F31; C2.5, C3.3, C3.4 \\
\end{longtable}
\endgroup

\subsection{Soluções descartadas}
\label{sub:r2-aprov-descartadas}

O estudo também registrou as soluções que a equipe decidiu não aproveitar, para que a decisão fique rastreável nas etapas seguintes do projeto. Quatro delas foram descartadas por evidência de uso ou de regulação. A visão computacional, base de parte das lojas sem caixa, não serve bem a roupas: a Amazon passou a usar RFID no Just Walk Out para vestuário porque a visão computacional exige produtos em prateleira \cite{amazon2024jwo}. O alarme sem barreira física também foi descartado, porque a Zara atrasou em 2023 a implantação de um sistema RFID sem etiqueta rígida, já que os chips eram fáceis de arrancar e o alarme tocava em devoluções \cite{pymnts2023zara}, e porque testes de retirada de etiquetas rígidas relatados pela Sensormatic, parte interessada, registraram aumento de perda de 40\% a mais que o dobro \cite{sensormatic2024hardtag}. O pagamento facial próprio, como o oferecido pela C\&A aos clientes do C\&A Pay \cite{mobiletime2023cea}, exigiria tratar dado biométrico, que a LGPD classifica como sensível (art.~11) \cite{brasil2018lgpd}; na carteira digital, a biometria fica no aparelho do cliente e não chega à Riachuelo. O Pix por aproximação, por fim, não se aplica, porque a etiqueta não é terminal de pagamento; o recurso começou em 2025, só no Android com Google Pay \cite{agenciabrasil2025pix}.

As outras cinco foram descartadas por limitação física, de escala ou de protótipo. A colheita de RF ambiente levaria cerca de 35\,h por liberação, segundo o balanço de energia dos estudos anteriores da equipe. Uma porta USB em cada etiqueta obrigaria a conectar à mão, uma a uma, as etiquetas da frota, o que não escala no recondicionamento. O anúncio BLE contínuo, a cada 1\,s, esgotaria a célula de 150\,mAh em cerca de 16 dias, estimativa da equipe, menos que um ciclo de 84 dias. O leitor UHF próprio no protótipo foi descartado porque o módulo avaliado, de US\$~79, não tem venda no Brasil (estudos anteriores da equipe); em escala, a leitura do EPC fica com o leitor da Estação de Cadastro e com a infraestrutura RFID que a Riachuelo já opera. O destacador magnético, por fim, foi substituído pelo carretel não ferromagnético e pelo Liberador de Balcão, porque desacopladores de 10.000\,GS são vendidos livremente por cerca de R\$~300 \cite{canalautomacao2026desacoplador} e abririam qualquer etiqueta com carretel ferromagnético.

\subsection{Propriedade intelectual}
\label{sub:r2-aprov-patentes}

A equipe fez a busca de patentes no Google Patents em 25/09/2026, a partir das referências das linhas L2 e L5, e conferiu a família de cada documento. O documento mais próximo do Tag\&Go é a patente US 10.121.339 B2, da Sensormatic (antiga Tyco), com prioridade em 04/03/2015, concedida em 06/11/2018 e ativa até 21/10/2035 \cite{googlepatents2018us10121339}. Sua reivindicação 1 descreve um celular, uma compra verificada, um comando sem fio e um pino liberado sem ajuda humana, conjunto muito próximo da concepção A1 Etiqueta Ativa. A família inclui os documentos CN107667395 (A e B), EP3265631 (A1 e B1), ES2746916T3, WO2016141039A1, US10121338B2 e US10522016B2, além de pedidos nos Estados Unidos, sem nenhuma publicação brasileira.

Outras três famílias ativas tratam da retirada de etiquetas depois do pagamento ou da energia na liberação. A US 10.354.507 B1, da Mastercard, descreve uma estação que remove a etiqueta depois de verificar o pagamento \cite{googlepatents2019us10354507}, e as US 11.417.186 B2 e US 11.984.003 B2, da Kohl's, descrevem a retirada automática de etiquetas \cite{googlepatents2022us11417186}; as duas famílias têm documentos só nos Estados Unidos. A US 10.497.238 B2, da Sensormatic, descreve a estação que energiza uma etiqueta sem bateria e tem família na China, na Europa, em Hong Kong e na Coreia do Sul, também sem publicação brasileira \cite{googlepatents2019us10497238}. Duas patentes da Checkpoint já expiraram: a US 7.450.013 B2, em 2025, hoje em domínio público e aproveitada na linha L2 \cite{googlepatents2008us7450013}, e a US 7.564.360 B2, que descreve um atuador dentro da etiqueta acionado por um sinal de liberação aberto, isto é, sem autenticação \cite{googlepatents2009us7564360}.

Em relação à US 10.121.339 B2, a solução do Tag\&Go difere em quatro pontos técnicos. O desafio (nonce) é gerado pela própria etiqueta, e o token é verificado nela mesma, com HMAC e chave por etiqueta, sem consulta à Plataforma Tag\&Go no momento da validação, de modo que o celular apenas retransmite e não carrega segredo. A associação entre etiqueta e EPC é feita na origem, e o pórtico RFID fecha a malha antifraude ao alarmar quando sai um EPC não vendido. A liberação assistida acontece por contatos, no Liberador de Balcão, mesmo com a célula descarregada. Por fim, os dois caminhos de uso, Compra Rápida e Jornada Completa, compartilham o mesmo núcleo comum. Essa comparação descreve diferenças de solução e não constitui parecer de liberdade de operação.

Como a proteção patentária é territorial, a equipe registrou apenas o que encontrou na consulta: nenhuma das famílias ativas tem publicação brasileira no Google Patents. Essa ausência ainda precisa ser confirmada no INPI, porque a fase nacional de um pedido internacional (PCT), como o WO2016141039A1 da família da US 10.121.339 B2, pode não aparecer no Google Patents. O efeito dessas famílias sobre fabricar e vender no Brasil, ou sobre exportar para os países em que elas existem, depende dessa confirmação e de uma análise jurídica que não faz parte deste relatório. A busca no INPI ficou como tarefa do Relatório 3, sob responsabilidade da dupla que elaborou este estudo. O risco tem precedente no varejo de moda: segundo os estudos anteriores da equipe, a Asterisk processou a Fast Retailing no Japão pela caixa de leitura RFID do self-checkout da Uniqlo e venceu a disputa em 2021.
